\documentclass[aspectratio=169,11pt]{beamer}
\usetheme{metropolis}

\usepackage{amsmath,amssymb}
\usepackage{booktabs}
\usepackage{graphicx}
\usepackage{xcolor}

\definecolor{mDarkTeal}{HTML}{1c3d6b}
\definecolor{mLightBrown}{HTML}{2f6fb3}
\definecolor{mLightGreen}{HTML}{4a90d9}
\setbeamercolor{structure}{fg=mDarkTeal}
\setbeamercolor{title}{fg=mDarkTeal}
\setbeamercolor{frametitle}{bg=mDarkTeal, fg=white}
\setbeamercolor{alerted text}{fg=mLightBrown}

\newcommand{\R}{\mathbb{R}}
\newcommand{\Gop}{\mathcal{G}}
\newcommand{\sgrad}{\operatorname{sg}}

\title{Improving Long-Rollout Accuracy of the DNO Surrogate}
\subtitle{Two training-time approaches we are evaluating}
\date{\today}
\author{Jonathan Ma}

\begin{document}

\begin{frame}\titlepage\end{frame}

% ============================================================
\begin{frame}{The rollout problem}

The surrogate learns to predict the Dirichlet--Neumann operator
$\Gop(\eta)\xi$ from snapshots of $(\eta,\xi,h)$, then is iterated
forward in time with a Zakharov integrator.

\medskip

Pointwise (single-step) accuracy is high. \alert{Long-time rollouts
are not.}

\pause
\medskip

On evaluation regimes (Tanaka solitons, Benjamin--Feir, Stokes,
random sea, linear), we observe error that grows
\emph{coherently in time}, dominated by modes the model is
\emph{supposed} to resolve --- not by aliasing, not by high-$k$
noise.

\medskip

The question: how do we train the surrogate to remain accurate when
fed its own predictions?
\end{frame}

% ============================================================
\begin{frame}{Why one-step training is insufficient}

Standard objective: given a data pair $(\eta,\xi) \mapsto \Gop(\eta)\xi$,
minimise a Sobolev-$H^k$ loss against the analytical target.

\medskip

\textbf{What this misses:}
\begin{itemize}
  \item The model only ever sees \alert{data-distribution} inputs.
  \item At rollout, it consumes its own prediction as input.
  \item Small biases at one step \alert{compound coherently} into the
        next step's input distribution.
\end{itemize}

\pause
\medskip

A loss that is locally accurate is not the same as a dynamical
system that is globally stable. The two approaches below directly
target this gap.
\end{frame}

% ============================================================
\begin{frame}{Approach 1: Pushforward training}

\textbf{Idea.}
Force the surrogate to be accurate at states \emph{it would itself
produce} after a few rollout steps.

\medskip

Given a data point $(\eta,\xi)$, advance the surrogate's own
prediction through $k$ Euler steps of the Zakharov system, then
re-predict the DNO at the rolled-out state:
\[
  (\eta_k,\xi_k) \;=\; \mathrm{Step}^k\bigl(\eta,\xi;\;
  \Gop_{\text{pred}}\bigr),
  \qquad
  \mathrm{pred}_k \;=\; \Gop_\theta\!\bigl(\sgrad(\eta_k),\sgrad(\xi_k)\bigr).
\]

The loss adds a term comparing $\mathrm{pred}_k$ against the analytical
Craig--Sulem expansion at the rolled-out state:
\[
  \mathcal{L} \;=\; \mathcal{L}_{\text{single-step}}
  \;+\; \lambda\,\mathcal{L}\!\bigl(\mathrm{pred}_k,\;
       \Gop_{\text{CS}}(\eta_k,\xi_k)\bigr).
\]

Gradients flow only through the final prediction: the trajectory and
the target are treated as fixed.

\end{frame}

% ============================================================
\begin{frame}{Approach 1: Pushforward --- intuition and limits}

\textbf{What it teaches.} The surrogate must be accurate not just on
the training manifold, but on the (slightly off-manifold) states it
itself generates. This breaks the covariate-shift compounding
explicitly.

\medskip

\textbf{What it does not address.}
\begin{itemize}
  \item No notion of physical conservation laws.
  \item Small $k$ ($k{=}1$) may not capture long-horizon coupling;
        large $k$ is expensive and harder to optimize.
  \item Stop-gradient on the rolled-out state introduces a known bias
        scaling like $k^2$ relative to the true multi-step gradient
        (List, Thuerey et al., 2024).
\end{itemize}

\pause
\medskip

\textbf{Status.} Currently trained at $k{=}1$, $\lambda{=}1$. A
longer-horizon, gradient-through unroll is a planned next step.
\end{frame}

% ============================================================
\begin{frame}{Approach 2: Hamiltonian-conservation loss}

\textbf{Idea.} The Zakharov system is Hamiltonian, with
\[
  H(\eta,\xi) \;=\; \tfrac12\!\int_\Omega
    \bigl(\xi\,\Gop(\eta)\xi \,+\, g\,\eta^2\bigr)\,dx.
\]
An exact solution conserves $H$. A surrogate whose rollout drifts in
$H$ is producing unphysical energy creation or destruction --- even
if its pointwise error per step is small.

\medskip

Add a loss term that penalises one-step relative drift of $H$:
\[
  \mathcal{L}_H \;=\; \lambda_H\,
    \mathbb{E}\!\left[\,\min\!\left(
      \frac{|H_1 - H_0|}{|H_0| + \varepsilon},\;C
    \right)\right],
\]
with $H_0$ evaluated at the data state and $H_1$ at one Euler step of
the surrogate's own dynamics. Gradient flows through both
$\Gop$-predictions; the integrator is treated as fixed.
\end{frame}

% ============================================================
\begin{frame}{Approach 2: H-cons --- intuition and limits}

\textbf{What it teaches.} A physically-motivated soft constraint:
``predictions that violate energy conservation are penalised, in
proportion to how much they violate it.''

\medskip

\textbf{What it does not address.}
\begin{itemize}
  \item Single-step horizon only --- does not directly stabilise
        long rollouts the way pushforward does.
  \item $H$ depends on $\Gop$, which is what we are learning ---
        so the loss is self-consistency of the surrogate's energy
        belief, not conservation of the true energy.
  \item Does not enforce momentum, mass, or other invariants.
\end{itemize}

\pause
\medskip

\textbf{Status.} Currently training with linear warmup and per-sample
clip to bound early-training transients. First evaluation pending.
\end{frame}

% ============================================================
\begin{frame}{What we see so far}

\textbf{Pushforward, vs.\ baseline at matched compute (partial).}
\begin{itemize}
  \item Median relative $L_2$ error at the final time horizon is
        reduced by roughly $30$--$70\%$ on Tanaka, Benjamin--Feir,
        and the linear regime.
  \item Stokes regimes are essentially unchanged (error already small).
  \item \alert{Random sea regimes regress} --- pushforward at $k{=}1$
        appears to trade noise-robustness on broadband ICs for
        fidelity on coherent, soliton-like ICs.
\end{itemize}

\pause
\medskip

\textbf{Caveats.}
\begin{itemize}
  \item Only $16$ ICs per regime --- per-regime comparisons have wide
        confidence intervals.
  \item Pushforward run is still training to matched epoch count.
  \item H-cons evaluation not yet available.
\end{itemize}
\end{frame}

% ============================================================
\begin{frame}{What we are working on next}

\begin{itemize}
  \item Finish the pushforward run to matched compute against the
        baseline.
  \item Evaluate the H-conservation run on the same suite.
  \item Investigate the random-sea regression: is it a
        pushforward-horizon issue, or does it need a different IC
        sampling distribution during training?
  \item Larger $k$, gradient-through unrolls (list, Thuerey 2024
        suggests $k\geq 8$ is where this becomes stable).
  \item Tendency-target reformulation (predict $\partial_t u$, then
        integrate with RK4 externally) --- decouples the surrogate
        from the dispersive phase rotation.
\end{itemize}
\end{frame}

\end{document}
